\section{Counting the multiplicative family}\label{sec:counting}

The first infinitude proof is immediate: take \(a=\ell\) prime with
\(\ell\equiv1\pmod9\), and invoke Dirichlet.  That already puts infinitely
many Hasse failures on one map.  But it throws away most parameters allowed
by the local calculation.  It proves infinitude, but it does not yet reveal
the natural size of the family.

The support condition in \eqref{eq:parameter-family} is multiplicative.
Products of suitable primes remain suitable, provided their residue modulo
\(9\) returns to \(1\), and the number of available prime factors is
unbounded.  The correct scale is therefore not the prime-counting scale
\(B/\log B\), but a Selberg--Delange scale.  We first translate parameter
size into target height and then count the semigroup.

We use the standard height induced by
\[
 \A^3(\Q)\longrightarrow\mathbb P^3(\Q),\qquad
 (u,v,w)\longmapsto[1:u:v:w].
 \tag{4.1}
\]
For \(a\in\mathcal A\), equation \eqref{eq:target-family} gives
\[
 [1:y_a]=[9:-9:32a:24a+3].
 \tag{4.2}
\]
Because \(a\equiv1\pmod9\), these coordinates are primitive; for \(a>1\),
their maximum absolute value is \(32a\).  Hence
\[
 H(y_a)=32a.
 \tag{4.3}\label{eq:height}
\]

\subsection{The full sufficient family}

Ignore the congruence and cube conditions for a moment, and count all
positive integers supported on primes \(1\pmod3\).  Their Dirichlet series
records the density \(1/2\) of the allowed primes:
\[
 \mathcal F_3(s)
 =\prod_{p\equiv1\ (3)}(1-p^{-s})^{-1}
 =\zeta(s)^{1/2}G_3(s).
 \tag{4.4}\label{eq:f3}
\]
If \(\chi_{-3}\) is the nonprincipal character modulo \(3\), comparison of
Euler factors gives
\[
 G_3(s)=
 \left[
 L(s,\chi_{-3})(1-3^{-s})
 \prod_{p\equiv2\ (3)}(1-p^{-2s})
 \right]^{1/2}.
 \tag{4.5}
\]
The branch is chosen positive for real \(s>1\).  In particular,
\[
 G_3(1)=
 \left[
 \frac23L(1,\chi_{-3})
 \prod_{p\equiv2\ (3)}(1-p^{-2})
 \right]^{1/2}>0.
 \tag{4.6}
\]

The Selberg--Delange theorem applied to \eqref{eq:f3} yields
\[
 \#\{n\le X:p\mid n\Rightarrow p\equiv1\pmod3\}
 \sim\frac{G_3(1)}{\sqrt\pi}\frac{X}{\sqrt{\log X}}.
 \tag{4.7}
\]
Such integers lie in the order-three subgroup
\[
 H=\{1,4,7\}\subset(\Z/9\Z)^\times.
\]
Here is the character step explicitly.  If \(\psi\) is a character of
\(H\), put
\[
 \mathcal F_{3,\psi}(s)
 =\prod_{p\equiv1\ (3)}
   (1-\psi(p)p^{-s})^{-1}.
\]
The prime number theorem in the three reduced classes \(1,4,7\pmod9\)
shows that its Selberg--Delange exponent is
\[
 \rho_\psi=\frac16\sum_{r\in H}\psi(r).
\]
Thus \(\rho_\psi=1/2\) for the trivial character and
\(\rho_\psi=0\) for either nontrivial character.  The
character-twisted Selberg--Delange theorem therefore gives
\[
 \sum_{\substack{n\le X\\p\mid n\Rightarrow p\equiv1\ (3)}}
 \psi(n)
 =o\left(\frac{X}{\sqrt{\log X}}\right)
\]
for both nontrivial twists.  Orthogonality on \(H\) now gives equal main
terms in its three residue classes.  Therefore
\[
 \#\left\{
 \begin{array}{@{}c@{}}
 n\le X:n\equiv1\pmod9,\\
 p\mid n\Rightarrow p\equiv1\pmod3
 \end{array}
 \right\}
 \sim
 \frac{G_3(1)}{3\sqrt\pi}\frac{X}{\sqrt{\log X}}.
 \tag{4.8}
\]
There are \(O(X^{1/3})\) perfect cubes up to \(X\), so excluding cubes and
the value \(1\) does not affect the main term.  We have proved
\[
 \#\{a\in\mathcal A:a\le X\}
 \sim
 \frac{G_3(1)}{3\sqrt\pi}\frac{X}{\sqrt{\log X}}.
 \tag{4.9}
\]
Substituting \(X=B/32\) by \eqref{eq:height} proves
Theorem~\ref{thm:count-intro}.

\subsection{A clean subfamily}

Before returning to the prime line, it is useful to isolate a simpler,
slightly thinner multiplicative subfamily
\[
 \mathcal A_9=
 \{a>1:a\notin\Q^3,\ p\mid a\Rightarrow p\equiv1\pmod9\}.
 \tag{4.10}
\]
It is obtained from the semigroup generated by primes \(1\pmod9\) by
removing the perfect cubes and \(1\), and it is contained in
\(\mathcal A\).  Its Dirichlet
series has the form
\[
 \mathcal F_9(s)
 =\prod_{p\equiv1\ (9)}(1-p^{-s})^{-1}
 =\zeta(s)^{1/6}G_9(s),
 \tag{4.11}
\]
where \(G_9\) is holomorphic and nonzero near \(s=1\).  Equivalently,
\[
 G_9(1)=
 \lim_{s\to1^+}(s-1)^{1/6}
 \prod_{p\equiv1\ (9)}(1-p^{-s})^{-1}>0.
 \tag{4.12}
\]
Another application of Selberg--Delange gives
\[
 \#\{a\in\mathcal A_9:a\le X\}
 \sim
 \frac{G_9(1)}{\Gamma(1/6)}
 \frac{X}{(\log X)^{5/6}}.
 \tag{4.13}
\]
The full family is larger because it may combine primes in all three
classes \(1,4,7\pmod9\), subject only to their product being \(1\pmod9\).

\subsection{Comparison with the prime subfamily}

Restricting to primes \(\ell\equiv1\pmod9\), the prime number theorem in
arithmetic progressions gives
\[
 \#\{\ell:H(y_\ell)\le B\}
 =\pi(B/32;9,1)
 \sim\frac{B}{192\log B}.
 \tag{4.14}
\]
Thus the multiplicative family improves the prime construction by a factor
of order \(\sqrt{\log B}\).  This improvement is arithmetic rather than
geometric: all targets still lie on the same rational line.  The line was
the geometric discovery; recognizing that its local criterion is
multiplicative is what changes the counting exponent.  The pleasure is that
the same local proof contains both facts once it is read in the right order.

The analytic input used above is the classical Selberg--Delange theorem in
the \(\zeta(s)^\rho G(s)\) form and its character-twisted version; we refer
to de la Bret\`eche and Tenenbaum \cite{BretecheTenenbaum}.
